% GENERATED FILE. Edit canonical lessons, then run npm run generate:books.
% canonical-source-hash: fnv1a64:4926b5a94396a251
% canonical-lessons: ZH-C383-dianfei, ZH-C383-langfei, ZH-C383-nengyuan, ZH-C383-jianshao, ZH-C383-yicixing

\chapter{Saving Energy}
\label{ch:zh-saving-energy}
\hlchaptermodality{383}

\begin{chapteropening}
\textbf{By the end of this chapter:} \emph{I can talk about the electricity bill, wasting things, energy, cutting down and throwaway things.}

Complete the last lesson of chapter 383: I can talk about the electricity bill, wasting things, energy, cutting down and throwaway things.
\end{chapteropening}

\section[diànfèi]{diànfèi — electricity bill}

\label{lesson:ZH-C383-dianfei}



\begin{quote}
Before the new one: say the Mandarin for to guarantee, then the Mandarin for to patrol.
\end{quote}

\subsection*{You'll want to know: diànfèi}
\emph{diànfèi} — \textquotedblleft{}electricity bill\textquotedblright{}.

\begin{quote}
\emph{xiàtiān kāi kōngtiáo diànfèi hěn guì} — In summer, with the air conditioning on, the electricity bill is high.
\end{quote}

Its characters come later; for now it is a word for the ear.

\begin{cousinweb}
\emph{diànfèi} is built from \zh{电} 'electricity' and \emph{fèi} 'a cost', both inherited from Middle Chinese, \emph{denH} and \emph{phj+jH} (Guangyun, in Baxter's transcription), which glosses the first 'lightning' and the second 'to use up'.
\end{cousinweb}

\begin{grammarlens}[title={What you've built}]
One more word for giving reasons.
\end{grammarlens}

\subsection*{Your turn}
\begin{itemize}
\raggedright
  \item \emph{Say it:} \emph{diànfèi}
  \item \emph{Say it:} \emph{diànfèi}, once more
  \item \emph{Say it:} \emph{xiàtiān kāi kōngtiáo diànfèi hěn guì} — In summer, with the air conditioning on, the electricity bill is high.
  \item \emph{Recall:} say the Mandarin for to guarantee, then the Mandarin for to patrol, then say \emph{diànfèi} again
\end{itemize}

\begin{tcolorbox}[breakable,colback=teal!4,colframe=teal!35!black,title={Before you move on}]
What does diànfèi mean? (\textquotedblleft{}Electricity bill\textquotedblright{}.) Say it once more.
\end{tcolorbox}
\section[làngfèi]{làngfèi — to waste}

\label{lesson:ZH-C383-langfei}



\begin{quote}
Before the new one: say the Mandarin for to patrol, then the Mandarin for the electricity bill.
\end{quote}

\subsection*{You'll want to know: làngfèi}
\emph{làngfèi} — \textquotedblleft{}to waste\textquotedblright{}.

\begin{quote}
\emph{xǐzǎo de shíhou bú yào làngfèi shuǐ} — Don't waste water when you have a shower.
\end{quote}

Its characters come later; for now it is a word for the ear.

\begin{cousinweb}
\emph{làngfèi} is built from \emph{làng} 'a wave' and \emph{fèi} 'to use up', both inherited from Middle Chinese, \emph{langH} and \emph{phj+jH} (Guangyun, in Baxter's transcription), which glosses the first 'a wave; reckless' and the second 'to spend down'.
\end{cousinweb}

\begin{grammarlens}[title={What you've built}]
One more word for giving reasons.
\end{grammarlens}

\subsection*{Your turn}
\begin{itemize}
\raggedright
  \item \emph{Say it:} \emph{làngfèi}
  \item \emph{Say it:} \emph{làngfèi}, once more
  \item \emph{Say it:} \emph{xǐzǎo de shíhou bú yào làngfèi shuǐ} — Don't waste water when you have a shower.
  \item \emph{Recall:} say the Mandarin for to patrol, then the Mandarin for the electricity bill, then say \emph{làngfèi} again
\end{itemize}

\begin{tcolorbox}[breakable,colback=teal!4,colframe=teal!35!black,title={Before you move on}]
What does làngfèi mean? (\textquotedblleft{}To waste\textquotedblright{}.) Say it once more.
\end{tcolorbox}
\section[néngyuán]{néngyuán — energy; sources of energy}

\label{lesson:ZH-C383-nengyuan}



\begin{quote}
Before the new one: say the Mandarin for the electricity bill, then the Mandarin for to waste.
\end{quote}

\subsection*{You'll want to know: néngyuán}
\emph{néngyuán} — \textquotedblleft{}energy; sources of energy\textquotedblright{}.

\begin{quote}
\emph{yīnwèi néngyuán hěn guì suǒyǐ diànfèi yě hěn guì} — Energy is expensive, so the electricity bill is high too.
\end{quote}

Its characters come later; for now it is a word for the ear.

\begin{cousinweb}
\emph{néngyuán} is built from \zh{能} 'able' and \emph{yuán} 'a source', both inherited from Middle Chinese, \emph{nong} and \emph{ngjwon} (Guangyun, in Baxter's transcription), which glosses the first 'skilled' and the second 'where water rises'. Energy as power for machines is a modern sense.
\end{cousinweb}

\begin{grammarlens}[title={What you've built}]
One more word for giving reasons.
\end{grammarlens}

\subsection*{Your turn}
\begin{itemize}
\raggedright
  \item \emph{Say it:} \emph{néngyuán}
  \item \emph{Say it:} \emph{néngyuán}, once more
  \item \emph{Say it:} \emph{yīnwèi néngyuán hěn guì suǒyǐ diànfèi yě hěn guì} — Energy is expensive, so the electricity bill is high too.
  \item \emph{Recall:} say the Mandarin for the electricity bill, then the Mandarin for to waste, then say \emph{néngyuán} again
\end{itemize}

\begin{tcolorbox}[breakable,colback=teal!4,colframe=teal!35!black,title={Before you move on}]
What does néngyuán mean? (\textquotedblleft{}Energy; sources of energy\textquotedblright{}.) Say it once more.
\end{tcolorbox}
\section[jiǎnshǎo]{jiǎnshǎo — to reduce; to cut down}

\label{lesson:ZH-C383-jianshao}



\begin{quote}
Before the new one: say the Mandarin for to waste, then the Mandarin for sources of energy.
\end{quote}

\subsection*{You'll want to know: jiǎnshǎo}
\emph{jiǎnshǎo} — \textquotedblleft{}to reduce; to cut down\textquotedblright{}.

\begin{quote}
\emph{zuò dìtiě kěyǐ jiǎnshǎo dǔchē} — Taking the underground can cut down on traffic jams.
\end{quote}

Its characters come later; for now it is a word for the ear.

Two third tones meet in this word: jiǎn is said with a rising tone before shǎo, and the spelling keeps both marks.

\begin{cousinweb}
\emph{jiǎnshǎo} is built from \emph{jiǎn} 'to lessen' and \zh{少} 'few', both inherited from Middle Chinese, \emph{keamX} and \emph{syewX} (Guangyun, in Baxter's transcription), which glosses the first 'to diminish' and the second 'not many'.
\end{cousinweb}

\begin{grammarlens}[title={What you've built}]
One more word for giving reasons.
\end{grammarlens}

\subsection*{Your turn}
\begin{itemize}
\raggedright
  \item \emph{Say it:} \emph{jiǎnshǎo}
  \item \emph{Say it:} \emph{jiǎnshǎo}, once more
  \item \emph{Say it:} \emph{zuò dìtiě kěyǐ jiǎnshǎo dǔchē} — Taking the underground can cut down on traffic jams.
  \item \emph{Recall:} say the Mandarin for to waste, then the Mandarin for sources of energy, then say \emph{jiǎnshǎo} again
\end{itemize}

\begin{tcolorbox}[breakable,colback=teal!4,colframe=teal!35!black,title={Before you move on}]
What does jiǎnshǎo mean? (\textquotedblleft{}To reduce; to cut down\textquotedblright{}.) Say it once more.
\end{tcolorbox}
\section[yícìxìng]{yícìxìng — disposable; used once}

\label{lesson:ZH-C383-yicixing}



\begin{quote}
Before the new one: say the Mandarin for sources of energy, then the Mandarin for to cut down on.
\end{quote}

\subsection*{You'll want to know: yícìxìng}
\emph{yícìxìng} — \textquotedblleft{}disposable; used once\textquotedblright{}.

\begin{quote}
\emph{wǒmen zuìhǎo bú yòng yícìxìng kuàizi} — We'd better not use disposable chopsticks.
\end{quote}

Its characters come later; for now it is a word for the ear.

Tone: \zh{一} is said yí here, rising, because cì after it is a falling tone.

\begin{cousinweb}
\emph{yícìxìng} is built from \zh{一} 'one', \emph{cì} 'a time' and \emph{xìng} 'nature', all inherited from Middle Chinese, \emph{'jit}, \emph{tshijH} and \emph{sjengH} (Guangyun, in Baxter's transcription), which glosses the second 'a stopping place' and 'an order'. The word is modern.
\end{cousinweb}

\begin{grammarlens}[title={What you've built}]
One more word for giving reasons.
\end{grammarlens}

\subsection*{Your turn}
\begin{itemize}
\raggedright
  \item \emph{Say it:} \emph{yícìxìng}
  \item \emph{Say it:} \emph{yícìxìng}, once more
  \item \emph{Say it:} \emph{wǒmen zuìhǎo bú yòng yícìxìng kuàizi} — We'd better not use disposable chopsticks.
  \item \emph{Recall:} say the Mandarin for sources of energy, then the Mandarin for to cut down on, then say \emph{yícìxìng} again
\end{itemize}

\begin{tcolorbox}[breakable,colback=teal!4,colframe=teal!35!black,title={Before you move on}]
What does yícìxìng mean? (\textquotedblleft{}Disposable; used once\textquotedblright{}.) Say it once more.
\end{tcolorbox}
